\bookchapter{第二十五章：中古スマートフォンの処理とデジタル遺産}

\chapterintro{ねえみんな、これまで20章以上の「修行」を経て、あなたは「デジタル初心者」からかなり落ち着いたスマートフォンユーザーへと成長したことでしょう。でも、テクノロジー生活には常に巡りがあります——新しい恋人もいつかは古い恋人になります。古い相棒を「引退」させようとするとき、どうやって堂々と、安全に別れを告げればいいか知っていますか？}

\booksection{一、古いスマートフォンの「安楽死」：安全な売却·譲渡ガイド}

核心原則：データ消去とは「削除」ではなく「上書き」です。写真をゴミ箱に入れてアルバムを空にすれば完了だと思う人が多いですが、スマホを譲渡した後、新しい持ち主がデータ復元ソフトでスキャンすると、古い写真、チャット履歴、身分証写真がすべて戻ってきます——このような「社会的死」現場は毎年上演されています。

\begin{enumerate}
  \item \textbf{一、古いスマートフォンの「安楽死」：安全な売却·譲渡ガイド}：最後の完全バックアップを完了し、すべてのアカウント（iCloud、Huawei/Xiaomi/Samsungアカウント、Googleアカウント）からサインアウトし、「スマホを探す」機能をオフにしてください——そうしないと次の使用者がアクティベーションロックで締め出されます。
  \item \textbf{自分で「大掃除」}：現代のスマホの「初期化」は通常、暗号化消去を含みます——暗号鍵自体を破体を破壊し、鍵がなくなると全データが瞬時に意味のない文字列になります。
  \item \textbf{最後の「目くらまし」（旧デバイスのみ必要）}：現代のスマホ（iOS 8+、Android 6+）はデフォルトでフルディスク暗号化（FBE）が有効で、初期化すると暗号鍵が破壊され、データは理論上復元不可能です。\textbf{この手順は不要}。極めて古い非暗号化デバイス（2015年以前）のみ、初期化後にジャンクファイルでストレージを満たし、再度初期化する必要があります。
\end{enumerate}

\textbf{各ブランドのデータ消去パス早見表}：原理は共通していますが、各ブランドで「初期化」のメニュー位置がわずかに異なります。主な違いは以下の通りです。

\begin{center}
\begin{tabularx}{\linewidth}{|X|X|X|}
\hline
\textbf{ブランド} & \textbf{消去エントリー} & \textbf{主な違いと注意事項} \tabularnewline
\hline
Huawei & 設定 > システムと更新 > リセット > 初期化 & 事前にHuaweiアカウントからサインアウトし「スマホを探す」をオフにする必要があります。そうしないとアクティベーションロックがトリガーされます。純正モードではパスワード認証が必要な場合があります。 \tabularnewline
Xiaomi / Redmi & 設定 > マイデバイス > 初期化 & 下部に「仮想SDカードをフォーマット」と「すべてのコンテンツと設定を消去」の両オプションがあり、両方ともチェックする必要があります。Xiaomiアカウントからサインアウトする前に「デバイスを探す」をオフにしてください。 \tabularnewline
OPPO / OnePlus & 設定 > その他の設定 > スマホをリセット > すべてのデータを完全に消去 & ColorOSは「完全消去」モードを提供し、ストレージを複数回上書きします。ロック画面のパスワード入力で確認が必要です。 \tabularnewline
vivo & 設定 > システム管理 > バックアップとリセット > すべてのデータを消去 & 「スマホのストレージをフォーマット」にチェックを入れるとSDカード領域も消去されます。vivoアカウントからサインアウトし、「スマホを探す」をオフにしてください。 \tabularnewline
Samsung & 設定 > 一般管理 > リセット > 初期化 & 事前にSamsungアカウントを削除し「スマホを探す」をオフにする必要があります。一部機種には「自動上書き」オプションがあり、有効化を推奨します。 \tabularnewline
iPhone & 設定 > 一般 > 転送またはiPhoneをリセット > すべてのコンテンツと設定を消去 & 事前にApple IDからサインアウト必須。そうしないと次のユーザーがiCloudアクティベーションロックで締め出されます。iOS 15+は「サインアウト時にiCloudデータを保持」に対応し、後で確認できます。 \tabularnewline
\hline
\end{tabularx}
\end{center}

\begin{tcbwarn}{最も見落とされやすい3つのこと}
\begin{itemize}
  \item \textbf{アカウントからサインアウトしていない}：スマホを買い手に売却した後に遠隔ロックされ、自分でも解除できず、最終的に身分証を持って公式サポートに行くしかなくなります。
  \item \textbf{SDカードを取り出していない}：初期化では外付けSDカード内の写真はフォーマットされません。手動で取り出すか、個別にフォーマットしてください。
  \item \textbf{USBメモリ/OTG外付けデバイスを切断していない}：消去中に外付けデバイスが接続されていると、外付けデバイスのデータが誤消去されたり、消去処理が中断される可能性があります。
\end{itemize}

\end{tcbwarn}

\booksection{二、デジタル遺産：あなたのソーシャルな魂をどう安置するか？}

スマホの中のデータは、もはや「仮想」のものではありません。あなたの資金、ソーシャル関係、創作物、さらには「デジタルアイデンティティ」を担っています。\textbf{人が亡くなった後、これらの資産をどう処理すべきでしょうか？}これはもはや抽象的な問題ではなく、ますます多くの家庭が現実に直面する難題です——家族がパスワードを知らず、プラットフォームが継承を許可せず、暗号資産が永遠に失われます。

\begin{tcbinfo}{法的根拠：個人情報保護法第49条}
2021年に施行された「中華人民共和国個人情報保護法」第49条は明確に規定しています：\textbf{自然人が死亡した場合、その近親者は自身の合法かつ正当な利益のために、死者の関連個人情報について本章に規定する閲覧、複製、訂正、削除等の権利を行使できる。ただし、死者が生前別の取り決めをした場合を除く。}

これは以下を意味します：

\begin{itemize}
  \item \textbf{近親者には法的根拠があり}プラットフォームに死者のアカウントと個人情報の処理を申請できます。
  \item 「合法かつ正当な利益」には通常、遺産相続の手続き、アカウントの閉鎖、機密プライバシーの削除などが含まれます。
  \item 死者が生前に遺言や契約で別の取り決め（遺産連絡先の指定など）をしていた場合、その取り決めに従います。
  \item プラットフォームは「ユーザープライバシー」を理由に一律に協力を拒否することはできませんが、具体的なフローは各社で異なります。
\end{itemize}

\end{tcbinfo}

\booksubsection{WeChatとAlipay：アカウントは継承できないが、資金は可能}

この2つの中国人が最もよく使うアプリは、ルールがそれぞれ異なります：

\begin{itemize}
  \item \textbf{WeChatアカウント}：「Tencent WeChatソフトウェア利用許諾およびサービス契約」に基づき、アカウントの所有権はTencentに帰属し、ユーザーは使用権のみを有します。\textbf{アカウント自体は継承も譲渡もできません}。ただしアカウント内の\textbf{残高、理財通の資金}は個人財産に属し、近親者は死亡証明、親族関係証明、相続公証書を用いてTenPayに引き出しまたは移転を申請できます。
  \item \textbf{Alipay}：アカウントは同様に直接継承できませんが、Alipay残高、Yu'ebao、理財通などの\textbf{資金は相続可能な財産に属します}。相続人はAlipayの「カスタマーサービスホットライン」または窓口で死亡証明、相続公証書を提出して手続きできます。Ant Fortune系の理財商品はファンド会社のフローを経る必要があります。
  \item \textbf{チャット履歴}：法的には個人のプライバシーに属し、プラットフォームが自発的に家族に提供することはありません。遺産の手がかりに関わる場合、訴訟や公証申請で取得できます。
\end{itemize}

\begin{tcbwarn}{見知らぬ人に「画面共有」や「リモートコントロール」を有効化。}
\begin{itemize}
  \item \textbf{死者のアカウントのパスワードを家族に直接渡して「そのまま使い続ける」のは避ける}——プラットフォーム契約違反となり、資金が凍結されたりリスク管理がトリガーされた際に立証がより困難になります。
  \item \textbf{死者のスマホで認証コードを受け取って「引き出す」のは避ける}——不正使用と認定され、かえって法律に抵触する可能性があります。
  \item \textbf{引き延ばさない}：一部のアカウントは長期間未ログインで自動的に停止され、資金の回収がより困難になります。
\end{itemize}

\end{tcbwarn}

\booksubsection{Apple「レガシーコンタクト」：事前に指定すれば労半して功倍す}

iOS 15.2以降、Appleは「レガシーコンタクト」機能を提供し、ユーザーが事前に信頼する連絡先を指定できます。亡くなった後、その連絡先は\textbf{死亡証明 + レガシーコンタクトアクセスキー}で死者のiCloudデータ（写真、メモ、連絡先など）へのアクセスを申請できます。操作手順：

\begin{enumerate}
  \item iPhoneで\textbf{[設定] > [ユーザー名] > [サインインとセキュリティ] > [レガシーコンタクト]}。
  \item 「レガシーコンタクトを追加」をタップし、Face IDで本人確認します。
  \item 連絡先を選択（連絡先から選択）すると、システムが\textbf{アクセスキー}。
  \item アクセスキーをメッセージ、メール、または印刷した紙面で相手に渡し、大切に保管するよう伝えてください。
  \item 複数の連絡先をバックアップとして追加でき、単一障害点を回避できます。
\end{enumerate}

指定された人はAppleデバイスがなくても、キーと死亡証明でApple公式サイトからアクセスを申請できます。全体のフローは通常7日以内に完了します。

\booksubsection{Google「非アクティブアカウントマネージャー」：亡き後の事務を自動処理}

Googleは「非アクティブアカウントマネージャー（Inactive Account Manager）」を提供し、ユーザーが長期間ログインしない場合に事前設定した操作を自動実行できます——これは現在最も完成されたデジタル遺産メカニズムの一つです。

\begin{enumerate}
  \item github.com/yourname\textbf{myaccount.google.com/inactive}、Googleアカウントにログインします。
  \item タイムアウト期間を設定：3 / 6 / 12 / 18ヶ月未ログインでトリガー。
  \item 信頼できる連絡先を追加（最大10人）、各人がアクセスできるデータ（Gmail、写真、Googleドライブなど）を設定。
  \item オプション：自動返信メールを設定し、差出人に対応できないことを通知。
  \item 最後に選択：期限到達後に\textbf{連絡先にデータを共有する}、それとも\textbf{アカウントを直接削除する}。
  \item 電話番号をフォールバック認証として設定し、誤トリガーを防ぎます。
\end{enumerate}

\booksubsection{国内ソーシャルアカウント：追悼化と削除フロー}

国内の主要ソーシャルプラットフォームは逝者のアカウント処理を2つのアプローチに分けます：\textbf{「追悼化」（アカウントを残し親族が偲べるようにする）}3. タスク（Task）：AIに何をするか明確に伝える\textbf{「削除」（完全に削除）}。家族は死亡証明、親族関係証明、申請人の身分証でプラットフォームのカスタマーサービスに申請できます。

\begin{itemize}
  \item \textbf{Weibo}：アカウントの「追悼化」に対応し、親族が死亡証明を提出して申請します。アカウントはロックされてログイン不可になりますが内容は保持され、プロフィールに「このアカウントは追悼アカウントに設定されています」と表示されます。削除申請も可能で、相続公証の提出が必要です。
  \item \textbf{Douyin}：「私 - 設定 - フィードバックとヘルプ - アカウントとセキュリティ - アカウント削除」またはカスタマーサービス経由で、死亡証明、親族関係証明、申請人の身分証を提出し、審査通過後に削除または追悼状態に移行できます。
  \item \textbf{Bilibili}：カスタマーサービスチャネルで逝者の証明と親族証明を提出し、アカウント削除または追悼ページとしての保留を申請できます。UP主の過去の動画の処理は著作権宣言と組み合わせる必要があります。
  \item \textbf{小紅書}：公式カスタマーサービスに申請し、死亡証明、親族関係証明を提出してアカウントの削除とコンテンツの削除を申請できます。
\end{itemize}

\begin{tcbinfo}{重要なリマインダー}
\begin{itemize}
  \item 各プラットフォームのルールは更新されるため、手続き前に公式カスタマーサービスの最新ガイダンスを確認してください。
  \item すべての証明書類は提出前にスキャンしてバックアップし、原本は大切に保管してください。
  \item 複数プラットフォームを一括手続きする際、統一的に\textbf{相続公証書}を1部申請すれば、手続きの往復を省けます。
\end{itemize}

\end{tcbinfo}

\booksubsection{デジタル資産の相続の難題}

ソーシャルアカウント以外にも、ますます多くの人が\textbf{真に価値のあるデジタル資産}を保有しており、その相続問題はより複雑です：

\begin{itemize}
  \item \textbf{暗号資産（Bitcoin、Ethereumなど）}：「秘密鍵が所有権」の原則に基づき、\textbf{秘密鍵がなければ資産は消滅したも同然}。相続人は法的文書でチェーン上に「申し立て」できません。助言：ニーモニック/秘密鍵を複数の物理的に安全な場所（銀行貸金庫+自宅金庫など）に分散保管し、遺言にアクセス方法を明記してください。
  \item \textbf{取引所アカウント}：銀行口座と同様に、相続公証書で取引所（Binance、OKXなど）に処分を申請できますが、越境取引所のフローは長引きます。
  \item \textbf{ドメイン名}：相続可能な財産に属し、相続人がドメイン登録業者（Alibaba Cloud、GoDaddyなど）に相続公証を提出して移転手続きができます。
  \item \textbf{知的財産権（オリジナル記事、動画、特許、ソフトウェア著作権）}：著作権の財産権は相続可能（人格権は非相続）で、収益権は相続人が著者の死後50年まで享受できます。
  \item \textbf{クラウドストレージとオンラインストレージ}：コンテンツデータは個人に帰属し、相続可能です。ただし一部のプラットフォーム契約ではアカウント自体の譲渡が不可とされ、先にダウンロードしてから削除する必要があります。
\end{itemize}

\textbf{暗号資産が「消滅」する真のコスト}\par

Chainalysisのレポートによると、世界中で既に\textbf{300万枚以上のBitcoin}が秘密鍵の紛失により永遠に回収不可能で、現存総量の6分の1に相当します。暗号資産を保有している場合、少なくとも1人の絶対的に信頼できる人に\textbf{復元パス}を知らせてください——さもなければこれらの富は永遠にブロックチェーン上で眠り続けます。

\booksubsection{「デジタル遺言」テンプレート}

従来の財産遺言ではデジタル資産を管理できません。毎年「デジタル遺言」添付書類を更新し、正式な遺言と一緒に保管することを推奨します。以下のテンプレートを参考にしてください：

\textbf{《個人デジタル資産処分添付書類》（テンプレート）}\par

\textbf{第1部：アカウントリスト}（アカウント名と所属プラットフォームのみ記録し、\textbf{パスワードは書かない}）

\begin{itemize}
  \item ソーシャル系：WeChat（アカウントxxx）、Weibo、Douyin、Bilibili、小紅書など
  \item 決済系：Alipay、WeChat Pay、各銀行カードアプリ
  \item クラウド系：iCloud、Huawei Cloud、Xiaomi Cloud、Googleアカウント、Microsoftアカウント
  \item メール系：QQメール、NetEaseメール、Gmailなど
  \item その他：オンラインストレージ、サブスクリプションサービス、ゲームアカウントなど
\end{itemize}

\textbf{第2部：パスワード管理方案}

\begin{itemize}
  \item マスターパスワードの保管方法：パスワードマネージャー（Bitwarden、1Passwordなど）を使用し、マスターパスワードは紙に書いて自宅の金庫に封印。
  \item 復旧キー/ニーモニック：分散保管し、少なくとも2カ所の物理的場所（銀行貸金庫+自宅金庫など）。
  \item 授権査問人：氏名、身分証番号、本人との関係、連絡先。
\end{itemize}

\textbf{第3部：処分の意思}

\begin{itemize}
  \item 各アカウントを「追悼化」するか「削除」するか。
  \item iCloud / Googleデータに「レガシーコンタクト」または「非アクティブアカウントマネージャー」を指定しているか。
  \item 暗号資産を誰が相続するか、秘密鍵の取得パスの説明。
  \item 著作権系作品の収益を誰が享受するか。
\end{itemize}

\textbf{第4部：連絡先とプラットフォーム}

\begin{itemize}
  \item 信頼する執行人の氏名、連絡先、予備連絡先。
  \item 常駐プラットフォームのカスタマーサービス窓口と必要な証明書類リスト。
  \item 本添付書類の更新日：\_\_\_\_年\_\_月\_\_日、次回確認日：\_\_\_\_年\_\_月\_\_日。
\end{itemize}

\begin{tcbwarn}{「パスワード」に関する法的注意}
パスワードを遺言に直接書き込むと漏洩リスクがあり、遺言公開後はすべての相続人が見られます。\textbf{より安全な方法は：遺言には「アカウントリスト + 処分の意思 + パスワード保管場所」のみを書き、実際のパスワードはパスワードマネージャーのマスターパスワードで信頼する執行人に別途伝える}。これにより法的要件を満たしつつ、パスワード漏洩リスクを低減できます。

\end{tcbwarn}

\booksection{三、エコな最後の旅路}

スマホ1台に金、銀、銅、希土類など数十種類の金属が含まれており、無断で廃棄すると土壌汚染と資源浪費につながります。正しいリサイクルルートは3種類あり、それぞれ特徴が異なります：

\begin{center}
\begin{tabularx}{\linewidth}{|X|X|X|X|}
\hline
\textbf{ルート} & \textbf{買取価格} & \textbf{データセキュリティ} & \textbf{：ハイエンドプロフェッショナル制作。iPhone 15 ProシリーズはApple ProRes動画撮影に対応（最大4K 60fps）しますが、ProRes RAW形式は現状外部録画デバイスとの組み合わせが必要で、スマホ単体での直接撮影は未対応。} \tabularnewline
\hline
公式下取り（Huawei/Apple/Xiaomiなど） & 中程度、補助券が併用されることが多い & 最高、公式の標準フローあり & 同ブランドの新機種に買い替える予定なら最もお得 \tabularnewline
サードパーティ回収プラットフォーム（愛回収、転転、回収宝など） & やや高い、価格比較可能 & 比較的高い、データ消去レポートあり & 直接現金化したい、ブランドに縛られたくない \tabularnewline
公益寄付（農村学校、慈善団体など） & 経済的リターンなし & 自分で完全消去が必要 & デバイスがまだ使える、善行をしたい \tabularnewline
再利用（自宅で改造） & 新デバイス費用を節約 & 元の所有者が使い続ける、リスクなし & デバイスの性能がまだ十分、家庭で改造ニーズがある \tabularnewline
\hline
\end{tabularx}
\end{center}

\begin{itemize}
  \item \textbf{：最も確実な方法、定期的に貴重なデータをローカルハードディスクにダウンロード。}：最も推奨される方法で、フローが標準化され、データセキュリティが保証されます。Apple、Huawei、Xiaomiなどのブランド公式サイトに「下取り」の入り口があり、評価後に新機種代に充当できます。
  \item \textbf{公式回収/下取り}：信用できる大型回収プラットフォームを選び、注文後に集荷に来てもらい、検査後に送金されます。必ず\textbf{「データ消去レポート」を提供する}プラットフォームを選び、証憑を保管してください。
  \item \textbf{公益寄付}：地元の農村学校やコミュニティ公益団体に連絡し、旧スマホを教育用タブレットや高齢者向けスマホに改造して使用できます。\textbf{寄付前に必ず完全な消去フローを実行し、すべてのSIM/SDカードを取り外してください}。
  \item \textbf{正規回収プラットフォーム}：車載ナビ、スマートホームハブ、監視カメラ、ベビーモニター、電子フォトフレームに改造。旧スマホ1台であと数年活用できます。
\end{itemize}

\begin{tcbinfo}{電子ゴミの真実}
スマホ1台の回路基板には\textbf{金、銀、パラジウム、銅}などの貴金属が含まれ、廃棄回路基板1トンあたりの金含有量は金鉱石の数十倍です。しかし無断廃棄すると、鉛、水銀、カドミウムが土壌と地下水に浸透し、汚染は数十年続きます。\textbf{正しくリサイクルすれば、環境を保護し、希少金属を循環再生できます。}

\end{tcbinfo}

\begin{tcbpractice}{活用}
\begin{enumerate}
  \item 最後の完全バックアップ：写真をクラウドにアップロード、重要ファイルをパソコンにコピー、チャット履歴をエクスポートして保存。
  \item すべてのアカウントからサインアウト：iCloud / Huawei / Xiaomi / Googleアカウント、「スマホを探す」をオフに。
  \item SIMカードとSDカードを取り出し、すべてのOTG外付けデバイスを切断。
  \item 暗号化初期化を実行、必要に応じて手動で「ストレージをフォーマット」にチェック。
  \item （旧デバイスのみ）ジャンクファイルでストレージを満たし、再度初期化。
  \item 次の行き先を選ぶ：公式下取り、正規回収プラットフォーム、公益寄付、または家族への譲渡。
  \item 「デジタル遺産リスト」を同期更新し、廃棄済みデバイスをリストから削除。
\end{enumerate}

\end{tcbpractice}

\begin{tcbpractice}{今月のタスクリスト}
\begin{enumerate}
  \item 現在のメインスマホにApple「レガシーコンタクト」を指定、またはGoogle「非アクティブアカウントマネージャー」を設定。
  \item 本節のテンプレートで《個人デジタル資産処分添付書類》を起草し、家族と保管場所を相談。
  \item 自宅に2台以上の未使用旧スマホがないか確認し、公式または正規プラットフォームの回収を予約。
\end{enumerate}

\end{tcbpractice}
